\section{Obligatory-testing instance optimality}
\label{sec:obligatory-instance}

The goal of this section is to prove \cref{thm:obligatory-instance}, the
rigorous version of the ``one algorithm for every input'' statement
\cref{thm:intro-obligatory-instance}.  The algorithm is exactly the
growing-cutoff algorithm $\mathcal A_n^*$ from
\cref{sec:random-obligatory}.  The new point is that, on every bounded input,
the algorithm attains the best leading cost available even to an announced
algorithm.

For a vector $p=(p_1,\ldots,p_n)$, we define below an explicit number
$\PhiOT(D[p])$, using the common empirical notation
\cref{eq:empirical-distributions}.  It is obtained by scanning the
processing times in increasing order, choosing the prefix with the largest
number of completions per unit of test-and-process work, and evaluating the
stationary algorithm for that prefix.

\begin{theorem}[Obligatory-testing instance optimality]
\label{thm:obligatory-instance}
There is a single randomized nonanticipating algorithm $\mathcal A^*$, depending
only on $n$ and not on the input vector, with the following property.  For every
fixed $L<\infty$ there is a deterministic sequence
$\eps_L(n)\to0$ such that every vector $p\in[0,L]^n$ satisfies
\begin{equation}
 \EE\ALG_{\mathcal A^*}(p)
 \le n^2\PhiOT(D[p])+\eps_L(n)n^2.
 \label{eq:ot-instance-upper}
\end{equation}
Conversely, every $p\in[0,L]^n$ and every randomized algorithm $\mathcal A$,
even an announced algorithm, satisfy
\begin{equation}
 \EE\ALG_{\mathcal A}(p)
 \ge n^2\PhiOT(D[p])-\eps_L(n)n^2.
 \label{eq:ot-instance-lower}
\end{equation}
One may take $\eps_L(n)=O_L(n^{-1/5})$.  Thus $\mathcal A^*$ is
asymptotically instance-optimal, uniformly over all bounded inputs.
\end{theorem}

Both estimates below are stated with the common order $O_L(n^{-1/5})$.

\begin{theorem}[When instance optimality is impossible]
\label{thm:ot-instance-impossibility}
\begin{enumerate}[label=\textnormal{(\roman*)}]
\item There is no deterministic instance-optimal algorithm, even when all job
lengths lie in $[0,2]$.  More precisely, for every deterministic algorithm
$\mathcal A$, even an announced one, and every even $n$, there is a fixed
labeling of the input with $n/2$ zero jobs and $n/2$ length-two jobs for which
\begin{equation}
 \ALG_{\mathcal A}\ge\frac98n^2.
 \label{eq:ot-det-not-instance-optimal}
\end{equation}
On the same input, an announced randomized algorithm $\mathcal A'$ has
\[
 \ALG_{\mathcal A'}=n^2+\frac34n.
\]
Thus $\ALG_{\mathcal A}>\ALG_{\mathcal A'}+n^2/8-O(n)$.

\item If job lengths are not bounded by a fixed constant, there is no
randomized instance-optimal algorithm with one uniform additive $o(n^2)$ term.
More precisely, for every randomized algorithm $\mathcal A$ independent of the
input and every sufficiently large $n$, there is an input with lengths in
$[0,n^2]$ and an announced randomized algorithm $\mathcal A'$
such that
\begin{equation}
 \ALG_{\mathcal A}\ge\ALG_{\mathcal A'}+\frac18n^2.
 \label{eq:ot-unbounded-not-instance-optimal}
\end{equation}
\end{enumerate}
\end{theorem}

\begin{proof}
As in the introduction, $\ALG_{\mathcal A}$ includes the expectation over
the private randomness of $\mathcal A$.  For part~(i), write $n=2m$.
Against the fixed deterministic algorithm, reveal
length two on the first $m$ labels it tests and zero on the remaining $m$.
After the interaction, record every answer on the label where it was given.
Rerunning the deterministic algorithm on this fixed labeling reproduces the
same actions and observations.
If the algorithm does not finish, the claim is immediate.  Otherwise, let $h$
be the number of length-two jobs completed before the first zero is tested.
By that time all $m$ positive jobs have been tested, so at least $m+2h$ work
has elapsed.  The first $h$ completed positive jobs contribute at least
$3h(h+1)/2$.  At that point the remaining jobs consist of $m$ unit tests of
zero jobs and $m-h$ processing operations of length two.  Even SPT order for
this remaining work gives
\begin{align*}
 \ALG_{\mathcal A}(p)
 &\ge \frac{3h(h+1)}2
 +(2m-h)(m+2h)+\frac{m(m+1)}2\\
 &\hspace{22mm}{}+m(m-h)+(m-h)(m-h+1)\\
 &=\frac92m^2+\frac12h^2+\frac32m+\frac12h
 \ge\frac98n^2.
\end{align*}
The comparison algorithm $\mathcal A'$ tests in a uniformly random order, completes zeros at
their tests, defers every length-two job, and finally processes the deferred
jobs.  This is the stationary algorithm with the zero jobs early.  Its exact
cost identity \cref{eq:stationary-prefix-cost} gives
\[
 \ALG_{\mathcal A'}=n^2+\frac34n.
\]
This proves the first claim.

For part~(ii), put $H=n^2$ and first run $\mathcal A$ on the input $p^0$ in
which every job has length $H$.  For a fixed private seed, let $X_k$ be the
number of completed jobs immediately before the $k$th test.  Every processing
operation contributes the same $Hn(n+1)/2$ in total, regardless of its
position.  The $k$th test delays $n-X_k$ unfinished jobs.  Hence the excess
over the algorithm $\mathcal A'$ that tests and processes each job
immediately satisfies
\begin{equation}
 \ALG_{\mathcal A}-\ALG_{\mathcal A'}
 =\EE\sum_{k=1}^n(k-1-X_k).
 \label{eq:ot-all-H-excess}
\end{equation}
If this sum has expectation at least $n^2/8$, we are done.

Otherwise, replace one uniformly random label by a zero and leave every
other length equal to $H$; call the resulting input $p^1$.  Conditional on the
seed, the rank $K$ at which the zero label is tested is uniform on
$\{1,\ldots,n\}$.  Until that test, the run agrees with the all-$H$ run, and
therefore $X_K$ length-$H$ jobs complete before the zero.  From
\cref{eq:ot-all-H-excess},
\begin{equation}
 \EE X_K
 =\frac1n\EE\sum_{k=1}^nX_k
 >\frac{n-1}{2}-\frac n8
 =\frac{3n}{8}-\frac12.
 \label{eq:ot-zero-rank-completions}
\end{equation}
Forcing the zero job behind $X_K$ length-$H$ jobs adds $HX_K$ to the
clairvoyant SPT lower bound.  On the other hand, an announced algorithm can test
in a private random order, defer positive jobs until it finds the zero, and
then finish the positive jobs.  Its completion order is SPT, and moving its
unit tests away from the corresponding processing operations costs at most
$n^2$.  Thus
\[
 \ALG_{\mathcal A}-\ALG_{\mathcal A'}
 \ge H\left(\frac{3n}{8}-\frac12\right)-n^2.
\]
For $H=n^2$ this exceeds $n^2/8$ for all sufficiently large $n$.  Averaging
over the possible zero labels is exactly the private-shuffle expectation
from the input convention.  This proves part~(ii).
\end{proof}

\subsection{The maximum-density benchmark}

We first define the benchmark for an arbitrary finitely supported
probability distribution $D$ on $[0,\infty)$.  Apply the maximum-density
module \cref{lem:maximum-density-module}, choose its threshold selector
$g_*$, and put
\begin{equation}
 a_*=a_D(g_*),\qquad w_*=w_D(g_*)=a_*\tau_D.
 \label{eq:ot-module-threshold}
\end{equation}
The residual measure $\nu$ from \cref{eq:common-module-data} has mass
$1-a_*$ and support in $[\tau_D,\infty)$.  For an empirical measure one
may equivalently sort the occurrences and maximize
\cref{eq:empirical-prefix-density} over processing-time prefixes $E$
containing all zeros.  A tied boundary class may be placed on either side,
by \cref{cor:empirical-density-prefix}.

Define
\begin{equation}
 \boxed{
 \PhiOT(D)
 =w_*\left(1-\frac{a_*}{2}\right)
  +\SPT(\nu).}
 \label{eq:ot-Phi}
\end{equation}
For $D=D[p]$ and a maximizing prefix $E$ made of whole job occurrences
rather than fractional mass, this is precisely $\operatorname{STAT}(E)$
from \cref{eq:stationary-leading-coefficient}.  The next
lemma proves both that the value is
independent of how probability mass at the threshold is split and that it
lower-bounds every fluid
algorithm.

\begin{lemma}[Pointwise completion envelope]
\label{lem:ot-completion-envelope}
For scaled time $s\ge0$, let
\begin{equation}
 \mathcal C_{\mathsf{OT},D}(s)=\max\left\{
  D(\{0\})T+\sum_{p>0}c_p:
  \begin{array}{l}
   0\le T\le1,\quad 0\le c_p\le D(\{p\})T,\\
   T+\sum_{p>0}p c_p\le s
  \end{array}\right\}.
 \label{eq:ot-envelope}
\end{equation}
Then $\mathcal C_{\mathsf{OT},D}$ is attained simultaneously at all times by
the following stationary schedule: repeatedly test an infinitesimal fraction
and immediately process the outcomes selected by $g_*$ until all jobs have
been tested, then process the residual measure $\nu$ in SPT order.
Moreover,
\begin{equation}
 \PhiOT(D)=\Area(\mathcal C_{\mathsf{OT},D}).
 \label{eq:ot-Phi-area}
\end{equation}
\end{lemma}

\begin{proof}
The envelope in \cref{eq:ot-envelope} is
$\mathcal C^v_{D,1}$ from \cref{eq:common-fluid-envelope}; the part
representing jobs handled without testing has zero capacity, so $v$ is
immaterial.  By
\cref{lem:test-module-envelope}, its item measure is
\[
                         a_*\delta_{\tau_D}+\nu.
\]
The optimal realization first performs the tests and selected processing
specified in the statement of the lemma and then processes $\nu$ in SPT
order, proving the pointwise assertion.  Since
$\nu$ has mass $1-a_*$ and is supported above $\tau_D$, the common area
formula \cref{eq:common-envelope-area} expands to
\begin{align*}
 \Area(\mathcal C_{\mathsf{OT},D})
 &=\tau_D\left(\frac{a_*^2}{2}+a_*(1-a_*)\right)
   +\SPT(\nu)\\
 &=w_*\left(1-\frac{a_*}{2}\right)
   +\SPT(\nu),
\end{align*}
because $w_*=a_*\tau_D$.  This is \cref{eq:ot-Phi} and is unaffected by
how probability mass at the boundary is split.
\end{proof}

\subsection{The announced upper bound}

Suppose the input is announced but its labeling is not.  Choose a maximizing
prefix $E$ made of whole job occurrences, test labels in a uniformly random
order, process a job immediately when it belongs to $E$, and finish the
remaining jobs in SPT order.  The exact
identity \cref{eq:stationary-prefix-cost} and
\cref{eq:ot-Phi} give
\begin{equation}
 \EE\ALG_{\mathrm{stat}}(p)
 =n^2\PhiOT(D[p])
  +\frac12\left(|E|+\sum_i p_i\right).
 \label{eq:ot-announced-upper}
\end{equation}
In particular, if $p\in[0,L]^n$, the remainder is at most
$(1+L)n/2$.  The announced optimum is therefore at most the claimed fluid
value plus $O_L(n)$.

\subsection{No adaptive announced algorithm has a better leading cost}

Uniformly shuffle the coordinates of $p$ and fix a
deterministic nonanticipating algorithm.  We may assume that it completes every
job, since otherwise the lower bound is immediate.  In the adaptive order in
which it tests new labels, the revealed occurrences form a uniform random
permutation $P_1,\ldots,P_n$ of the coordinates of $p$ by deferred decisions.
When values are equal, we still treat the corresponding jobs as distinct.
Because the algorithm first touches every job by testing it, no final part of
the schedule needs to be modified, and we may directly use the sampling bound
that holds simultaneously for every prefix.

Keep zero separate, partition $(0,L]$ into
\begin{equation}
 K_n=\lfloor n^{1/5}\rfloor
 \label{eq:ot-lower-grid-size}
\end{equation}
cells, and round every positive value upward to its cell endpoint.  Let
$D^+$ be the resulting empirical law and $h_n=L/K_n$ the mesh.  For each
cell $j$, let $H_{k,j}$ count its occurrences among the first $k$ tests and
let $d_j$ be its mass in $D^+$.  The uniform-prefix bound
\cref{eq:uniform-prefix-hoeffding} gives, except on an event of probability
at most $n^{-2}$,
\begin{equation}
 |H_{k,j}-kd_j|\le\Gamma_n
 \quad\text{for all }k,j,
 \qquad
 \Gamma_n=O\!\left(\sqrt{n\ln(n+2)}\right).
 \label{eq:ot-cell-discrepancy}
\end{equation}

Put $\gamma_n=\Gamma_n/n$.  At any point in the schedule, let $t=k/n$, let $z$ be
the normalized number of completed zero jobs, and let $c_j$ be the completed
positive mass from cell $j$.  Equation~\eqref{eq:ot-cell-discrepancy} gives
\[
 z\le d_0t+\gamma_n,\qquad c_j\le d_jt+\gamma_n.
\]
Set $c_j'=(c_j-\gamma_n)^+$.  Rounding upward adds at most $h_n$ to the
normalized processing work accumulated in any prefix, so $t,c_j'$ satisfy
the constraints defining $\mathcal C_{\mathsf{OT},D^+}$ at work
$x+h_n$.  Replacing $z,c_j$ by $d_0t,c_j'$ removes at most
$(K_n+1)\gamma_n$ completed mass.  Hence, simultaneously at every prefix,
\begin{equation}
 \frac{N_{\mathcal A}(nx)}n
 \le \mathcal C_{\mathsf{OT},D^+}(x+h_n)
      +(K_n+1)\gamma_n.
 \label{eq:ot-pointwise-completion}
\end{equation}

The normalized completion cost is the area under the unfinished fraction.
The completion curve reaches one by normalized work $H=L+2$, so integrating
\cref{eq:ot-pointwise-completion} loses at most
$h_n+H(K_n+1)\gamma_n$ from $\PhiOT(D^+)$.  The good event fails with
probability at most $n^{-2}$, and the immediate test--process algorithm gives
$0\le\PhiOT(D^+)\le(L+1)/2$.  Finally, zero-preserving transport
\cref{lem:zero-preserving-transport} changes the fluid value by at most
$h_n$.  Because
$(K_n+1)\gamma_n=O(n^{-3/10}\sqrt{\ln(n+2)})$, every deterministic announced
algorithm has private-shuffle expected cost at least
\begin{equation}
 n^2\PhiOT(D[p])-O_L(n^{9/5}).
 \label{eq:ot-announced-lower}
\end{equation}
The argument was uniform after fixing the algorithm seed.  Averaging over
that seed gives precisely the randomized cost under the input convention.

\subsection{Removing the announcement}

It remains to prove \cref{eq:ot-instance-upper}.  Use the same sampling,
threshold, and fallback rules as $\mathcal A_n^*$ in
\cref{sec:random-obligatory}.  Its full parameter choice is
\begin{equation}
 k=\lfloor n^{3/4}\rfloor,\qquad
 d=\lfloor n^{1/4}\rfloor,\qquad
 B_n=32+n^{1/20},\qquad \eta=B_n/d.
 \label{eq:ot-learning-parameters}
\end{equation}
The algorithm privately permutes the labels, tests the first $k$, and rounds
positive values upward to the $B_n$-grid.  It starts with a sample prefix that
maximizes density and then includes every finite grid category up to the
resulting threshold.

Fix $L$.  For all sufficiently large $n$, $B_n>L$ and the grid has no
value in its overflow category.  Moreover, the prefix containing the entire
sample has density at least
$1/(1+L+\eta)>1/B_n$, so fallback never occurs.  Let $D^+$ be the rounded
population histogram, $D^{\rm rem}$ the rounded histogram of the remaining jobs,
$\widehat D$ the sample histogram, and
$\Delta=\|D^+-\widehat D\|_1$.  By \cref{lem:rand-histogram} and deletion of
the $k$ sample occurrences,
\begin{equation}
 \EE\Delta\le\sqrt{\frac{d+2}{k}},\qquad
 \|D^{\rm rem}-D^+\|_1\le\frac{2k}{n-k}.
 \label{eq:ot-learning-histogram}
\end{equation}

For a fixed grid prefix $J$, let $F_{\mathsf{OT},D}(J)$ be the right-hand side of
\cref{eq:ot-Phi} when $J$ is used as the early set, whether or not it is
optimal for $D$.  If $D,E$ are grid histograms supported on $[0,L+\eta]$,
then
\begin{equation}
 |F_{\mathsf{OT},D}(J)-F_{\mathsf{OT},E}(J)|
 \le3(L+2)\|D-E\|_1.
 \label{eq:ot-template-lipschitz}
\end{equation}
Indeed, early mass changes by at most the $\ell_1$ distance, early first
moment by at most $(L+\eta)$ times that distance.  In the remaining double
sum of $\min\{p,q\}$, change the first distribution and then the second; the
total change is at most twice the same amount.  The displayed loose constant
covers the three terms.

By \cref{lem:ot-completion-envelope}, a maximum-density prefix minimizes
$F_{\mathsf{OT},D}(J)$ over grid prefixes.  Let $\widehat J$ be the prefix
learned from the sample.
Rounding each processing time upward by at most $\eta$ preserves zeros, so
\cref{lem:zero-preserving-transport} gives
\begin{equation}
 \PhiOT(D^+)\le\PhiOT(D[p])+\eta.
 \label{eq:ot-rounding-stability}
\end{equation}

Conditional on the unordered sample, compare the actual schedule, which
tests the sample before the other jobs, with the fully stationary algorithm
for $\widehat J$.  The calculation
in \cref{sec:random-obligatory} bounds the extra cost by
$O_L(nk+k^2)$.  The one-job contribution in the exact stationary identity is
$O_L(n)$, and using the rounded values can only increase the processing
delays in that stationary schedule.  Therefore
\begin{equation}
 \EE[\ALG_{\mathcal A^*}(p)\mid S]
 \le n^2F_{\mathsf{OT},D^{\rm rem}}(\widehat J)+O_L(nk+k^2+n).
 \label{eq:ot-sample-first-instance}
\end{equation}
This is the conditional hypothesis of
\cref{lem:pilot-learning-transfer}.  Apply that lemma with metric
$d(D,E)=\|D-E\|_1$, $L_0=3(L+2)$, and
\cref{eq:ot-learning-histogram}; then use
\cref{eq:ot-rounding-stability}.
Since
\[
 \sqrt{d/k}=O(n^{-1/4}),\qquad
 \eta=O(n^{-1/5}),\qquad k/n=O(n^{-1/4}),
\]
we obtain
\begin{equation}
 \EE\ALG_{\mathcal A^*}(p)
 \le n^2\PhiOT(D[p])+O_L(n^{9/5})
 \label{eq:ot-instance-upper-rate}
\end{equation}
for every $p\in[0,L]^n$.  Together with
\cref{eq:ot-announced-lower}, this proves
\cref{thm:obligatory-instance}.

Finally, the exact identity \cref{eq:rand-certificate} says that
$\PhiOT(D)$ is at most $4/3$ times the leading clairvoyant SPT value
for every $D$.  Equality is attained by
$D=\tfrac12\delta_0+\tfrac12\delta_2$.  Thus the sharp $4/3$ theorem follows
from the stronger instance-optimal description by maximizing over inputs.
